\chapter{Практическое задание 5. Нагрев шара постоянным потоком}

\section{Постановка и знак теплового потока}

Однородный шар радиуса $R$ имеет нулевую начальную температуру.
На его поверхность с момента $t=0$ равномерно подаётся постоянный
поток тепла. Обозначим через $q>0$ плотность входящего теплового потока
(теплота на единицу площади за единицу времени), через $k>0$ ---
теплопроводность, через $c>0$ --- объёмную теплоёмкость.
Положим $a^2=k/c$.
Если в условии $q$ означает полный поток через всю поверхность,
в формулах вместо $q$ следует использовать $q/(4\pi R^2)$.

Из сферической симметрии температура зависит только от радиуса $r$ и
времени. Уравнение и условия имеют вид
\[
\boxed{\begin{cases}
 u_t=a^2\left(u_{rr}+\dfrac2r u_r\right),&0<r<R,\ t>0,\\
 u_r(0,t)=0,\quad u(0,t)\text{ конечна},\\
 k u_r(R,t)=q,\\
 u(r,0)=0.
\end{cases}}
\]
Наружная плотность потока по закону Фурье равна $-ku_r(R,t)$.
Поскольку шар нагревается, наружный поток равен $-q$, откуда
$ku_r(R,t)=q$. Эта проверка определяет знак граничного условия.

\section{Баланс тепла и выделение растущей части}

Средняя по объёму температура равна
\[
 \overline u(t)=\frac3{R^3}\int_0^R u(r,t)r^2\,dr.
\]
Интегрирование уравнения даёт
\[
 \overline u'(t)=\frac{3a^2}{R^3}[r^2u_r]_0^R
 =\frac{3q}{cR},\qquad
 \boxed{\overline u(t)=\frac{3q}{cR}t.}
\]
Следовательно, при $q>0$ стационарная температура невозможна:
тепло непрерывно поступает в шар и не отводится.

Выделим точное частное решение
\[
 p(r,t)=\frac{3q}{cR}t+
 \frac{qR}{k}\left(\frac{r^2}{2R^2}-\frac3{10}\right).
\]
Оно удовлетворяет уравнению и условию $kp_r(R,t)=q$.
Постоянная $-3/10$ выбрана так, чтобы среднее пространственной
части было нулевым. Положим $u=p+v$. Тогда
\[
 v_t=a^2\left(v_{rr}+\frac2r v_r\right),\qquad
 v_r(0,t)=v_r(R,t)=0,
\]
\[
 v(r,0)=-\frac{qR}{k}\left(\frac{r^2}{2R^2}-\frac3{10}\right).
\]
Начальное среднее $v$ равно нулю; постоянная собственная мода отсутствует.

\section{Собственные функции и коэффициенты}

Разделение переменных приводит к задаче
\[
 (r^2X')'+\lambda r^2X=0,\qquad X\text{ регулярна при }r=0,
 \qquad X'(R)=0.
\]
Для положительных собственных значений удобно взять
\[
 X_n(r)=\frac{\sin(z_nr/R)}{z_nr/R},\qquad X_n(0)=1,
 \qquad\lambda_n=\frac{z_n^2}{R^2}.
\]
Условие при $r=R$ даёт
\[
 z_n\cos z_n-\sin z_n=0,
 \qquad \boxed{\tan z_n=z_n.}
\]
Здесь $z_n$ --- строго положительные ненулевые корни:
\[
 z_1\approx4.493409,\qquad z_2\approx7.725252,\qquad
 z_3\approx10.904122.
\]
Корень $z=0$ соответствует постоянной моде и рассматривается отдельно.
Ортогональность осуществляется с весом $r^2$.
Поэтому
\[
 v(r,t)=\sum_{n=1}^{\infty}D_nX_n(r)e^{-a^2z_n^2t/R^2},
 \qquad
 D_n=\frac{\int_0^R v(r,0)X_n(r)r^2\,dr}
 {\int_0^RX_n^2(r)r^2\,dr}.
\]
Положив $s=r/R$ и используя $\sin z_n=z_n\cos z_n$, получаем
\[
 \int_0^1s^2\frac{\sin(z_ns)}{z_ns}\,ds=0,
\]
\[
 \int_0^1s^2\left(\frac{s^2}2-\frac3{10}\right)
 \frac{\sin(z_ns)}{z_ns}\,ds=\frac{\cos z_n}{z_n^2},
\]
\[
 \int_0^1s^2\left(\frac{\sin(z_ns)}{z_ns}\right)^2ds
 =\frac{\cos^2z_n}{2}.
\]
Следовательно,
\[
 \boxed{D_n=-\frac{2qR}{kz_n^2\cos z_n}.}
\]

\section{Распределение температуры и асимптотика}

Для $0\le r\le R$, $t>0$ окончательно получаем
\[
\boxed{\begin{aligned}
 u(r,t)={}&\frac{3q}{cR}t+
 \frac{qR}{k}\left(\frac{r^2}{2R^2}-\frac3{10}\right)\\
 &-\frac{2qR}{k}\sum_{n=1}^{\infty}
 \frac{e^{-a^2z_n^2t/R^2}}{z_n^2\cos z_n}
 \frac{\sin(z_nr/R)}{z_nr/R}.
\end{aligned}}
\]
При $r=0$ отношение синуса к аргументу заменяется его пределом $1$:
\[
 u(0,t)=\frac{3q}{cR}t-\frac{3qR}{10k}
 -\frac{2qR}{k}\sum_{n=1}^{\infty}
 \frac{e^{-a^2z_n^2t/R^2}}{z_n^2\cos z_n}.
\]
На поверхности шара
\[
 u(R,t)=\frac{3q}{cR}t+\frac{qR}{5k}
 -\frac{2qR}{k}\sum_{n=1}^{\infty}\frac{e^{-a^2z_n^2t/R^2}}{z_n^2}.
\]
Ряд восстанавливает нулевую начальную температуру.
Включение ненулевого потока при $t=0$ нарушает согласование
начальных и граничных производных; граничное условие выполняется
для $t>0$.

При больших временах
\[
 u(r,t)=\frac{3q}{cR}t+
 \frac{qR}{k}\left(\frac{r^2}{2R^2}-\frac3{10}\right)
 +O\left(e^{-a^2z_1^2t/R^2}\right).
\]
Температура неограниченно растёт, но её отличие от средней стремится
к постоянному пространственному профилю. В частности,
\[
 u(R,t)-u(0,t)\longrightarrow\frac{qR}{2k}.
\]
Энергетическая проверка решения:
\[
 c\frac{4\pi R^3}{3}\overline u(t)=4\pi R^2qt.
\]
Левая часть равна накопленной теплоте, правая --- теплоте,
поступившей через поверхность шара.
